%%% Hlavní soubor. Zde se definují základní parametry a odkazuje se na ostatní části. %%%

%% Verze pro jednostranný tisk:
% Okraje: levý 40mm, pravý 25mm, horní a dolní 25mm
% (ale pozor, LaTeX si sám přidává 1in)
\documentclass[12pt,a4paper]{report}
\setlength\textwidth{145mm}
\setlength\textheight{247mm}
\setlength\oddsidemargin{15mm}
\setlength\evensidemargin{15mm}
\setlength\topmargin{0mm}
\setlength\headsep{0mm}
\setlength\headheight{0mm}
% \openright zařídí, aby následující text začínal na pravé straně knihy
\let\openright=\clearpage

%% Pokud tiskneme oboustranně:
% \documentclass[12pt,a4paper,twoside,openright]{report}
% \setlength\textwidth{145mm}
% \setlength\textheight{247mm}
% \setlength\oddsidemargin{15mm}
% \setlength\evensidemargin{0mm}
% \setlength\topmargin{0mm}
% \setlength\headsep{0mm}
% \setlength\headheight{0mm}
% \let\openright=\cleardoublepage

%% Použité kódování znaků: obvykle latin2, cp1250 nebo utf8:
\usepackage[utf8]{inputenc}

%% Ostatní balíčky
\usepackage{graphicx}
\usepackage{amsthm}
\usepackage{nath}

%% Balíček hyperref, kterým jdou vyrábět klikací odkazy v PDF,
%% ale hlavně ho používáme k uložení metadat do PDF (včetně obsahu).
%% POZOR, nezapomeňte vyplnit jméno práce a autora.
\usepackage[ps2pdf,unicode]{hyperref}   % Musí být za všemi ostatními balíčky
\hypersetup{pdftitle=Practical data structures}
\hypersetup{pdfauthor=Michael Pokorný}

%%% Drobné úpravy stylu

% Tato makra přesvědčují mírně ošklivým trikem LaTeX, aby hlavičky kapitol
% sázel příčetněji a nevynechával nad nimi spoustu místa. Směle ignorujte.
\makeatletter
\def\@makechapterhead#1{
  {\parindent \z@ \raggedright \normalfont
   \Huge\bfseries \thechapter. #1
   \par\nobreak
   \vskip 20\p@
}}
\def\@makeschapterhead#1{
  {\parindent \z@ \raggedright \normalfont
   \Huge\bfseries #1
   \par\nobreak
   \vskip 20\p@
}}
\makeatother

% Toto makro definuje kapitolu, která není očíslovaná, ale je uvedena v obsahu.
\def\chapwithtoc#1{
\chapter*{#1}
\addcontentsline{toc}{chapter}{#1}
}

\newtheorem{theorem}{Theorem}
\newtheorem*{define}{Definition}	% Definice nečíslujeme, proto "*"

\def\O{
\mathcal{O}
}
%\newcommand{O}[0]{\mathcal{O}}

\begin{document}

% Trochu volnější nastavení dělení slov, než je default.
\lefthyphenmin=2
\righthyphenmin=2

%%% Titulní strana práce

\pagestyle{empty}
\begin{center}

\large

Charles University in Prague

\medskip

Faculty of Mathematics and Physics

\vfill

{\bf\Large BACHELOR THESIS}

\vfill

\centerline{\mbox{\includegraphics[width=60mm]{img/logo.eps}}}

\vfill
\vspace{5mm}

{\LARGE Michael Pokorný}

\vspace{15mm}

% Název práce přesně podle zadání
{\LARGE\bfseries Practical data structures}

\vfill

% Název katedry nebo ústavu, kde byla práce oficiálně zadána
% (dle Organizační struktury MFF UK)
Department of Applied Mathematics

\vfill

\begin{tabular}{rl}

Supervisor of the bachelor thesis: & Mgr. Martin Mareš, Ph.D. \\
\noalign{\vspace{2mm}}
Study programme: & Computer science \\
\noalign{\vspace{2mm}}
Specialization: & General computer science \\
\end{tabular}

\vfill

% Zde doplňte rok
Prague 2015

\end{center}

\newpage

%%% Následuje vevázaný list -- kopie podepsaného "Zadání bakalářské práce".
%%% Toto zadání NENÍ součástí elektronické verze práce, nescanovat.

%%% Na tomto místě mohou být napsána případná poděkování (vedoucímu práce,
%%% konzultantovi, tomu, kdo zapůjčil software, literaturu apod.)

\openright

\noindent
Dedication. TODO

\newpage

%%% Strana s čestným prohlášením k bakalářské práci

\vglue 0pt plus 1fill

\noindent
I declare that I carried out this bachelor thesis independently, and only with the cited
sources, literature and other professional sources.

\medskip\noindent
I understand that my work relates to the rights and obligations under the Act No.
121/2000 Coll., the Copyright Act, as amended, in particular the fact that the Charles
University in Prague has the right to conclude a license agreement on the use of this
work as a school work pursuant to Section 60 paragraph 1 of the Copyright Act.

\vspace{10mm}

\hbox{\hbox to 0.5\hsize{%
TODO
In ........ date ............
\hss}\hbox to 0.5\hsize{%
TODO
signature of the author
\hss}}

\vspace{20mm}
\newpage

%%% Povinná informační strana bakalářské práce

\vbox to 0.5\vsize{
\setlength\parindent{0mm}
\setlength\parskip{5mm}

Název práce:
Praktické datové struktury
% přesně dle zadání

Autor:
Michael Pokorný

Katedra:  % Případně Ústav:
Katedra aplikované matematiky
% dle Organizační struktury MFF UK

Vedoucí bakalářské práce:
Mgr. Martin Mareš, Ph.D.
% dle Organizační struktury MFF UK, případně plný název pracoviště mimo MFF UK

Abstrakt:
% abstrakt v rozsahu 80-200 slov; nejedná se však o opis zadání bakalářské práce
TODO

Klíčová slova:
% 3 až 5 klíčových slov
TODO

\vss}\nobreak\vbox to 0.49\vsize{
\setlength\parindent{0mm}
\setlength\parskip{5mm}

Title:
Practical data structures
% přesný překlad názvu práce v angličtině

Author:
Michael Pokorný

Department:
Department of Applied Mathematics
% Název katedry či ústavu, kde byla práce oficiálně zadána
% dle Organizační struktury MFF UK v angličtině

Supervisor:
% dle Organizační struktury MFF UK, případně plný název pracoviště
Mgr. Martin Mareš, Ph.D.
% mimo MFF UK v angličtině

Abstract:
% abstrakt v rozsahu 80-200 slov v angličtině; nejedná se však o překlad
% zadání bakalářské práce
TODO

Keywords:
% 3 až 5 klíčových slov v angličtině
TODO

\vss}

\newpage

%%% Strana s automaticky generovaným obsahem bakalářské práce. U matematických
%%% prací je přípustné, aby seznam tabulek a zkratek, existují-li, byl umístěn
%%% na začátku práce, místo na jejím konci.

\openright
\pagestyle{plain}
\setcounter{page}{1}
\tableofcontents

%%% Jednotlivé kapitoly práce jsou pro přehlednost uloženy v samostatných souborech
\chapter*{Introduction}
\addcontentsline{toc}{chapter}{Introduction}

It is a well-known fact that there is a widening gap between
the performance of CPUs and memory. As an optimization of memory accesses,
modern computers include several small and fast caches between the CPU and
main memory.
However, if programs access data mostly randomly, I/O can become
the performance bottleneck, even for programs running in main memory.

Standard models of computation only measure the performance of algorithms
in the number of CPU instructions executed and the number of memory cells used.
These models aren't well suited for environments where different memory cells
may have vastly different access speeds, for example depending on which caches
currently hold the memory cell. The \textit{external memory model} is a simple
extension of the RAM model that add a \textit{block transfer} measure, which
is the number of \textit{blocks} transferred between a fast cache and a slow
\textit{external memory}. While this model was originally developed
in the context of accessing data stored or hard disks, it is also useful
for reasoning about data transfers between the CPU and main memory.

\textit{Cache-oblivious algorithms} are a class of external memory algorithms
that perform equally well given any cache size and block size. Thanks to this
property, they need little tuning for efficient implementation.
Modern CPU architectures perform various optimizations that can change
the effective size of a block, so tuning an algorithm by setting a good block
size might be difficult.

This thesis explores the performance implications of memory hierarchies
on data structures implementing dictionaries or sorted dictionaries.
Especially unsorted dictionaries are very common in practice, as illustrated
by modern programming languages, which commonly have a dedicated built-in
data type for unsorted dictionaries (e.g. \texttt{std::unordered\_map} in C++,
\texttt{dict} in Python, or hashes in Perl).

TODO: all logs are $\log_2$

\chapter{Title of the first chapter}

\section{Title of the first subchapter of the first chapter}

\section{Title of the second subchapter of the first chapter}


\chapter*{Conclusion}
\addcontentsline{toc}{chapter}{Conclusion}

We confirmed that cache-oblivious B-trees are practically competitive with
more common data structures in main memory, especially on larger dictionaries
with frequent \textsc{Find}s and uniform access patterns.

Our implementation of $k$-splay trees and $k$-forests performed badly.
In our opinion, they alone are not a practical replacement for simple splay
trees. A more optimized implementation of $k$-splaying might reach parity with
splay trees.

Experiments on data recorded from Mozilla Firefox suggest that splay trees
and cuckoo hash tables do well on most dictionaries as they are practically
used. A particularly interesting conclusion that can be drawn from
Table~\ref{tab:firefox-results} is that hashing with linear probing may
be (somewhat counterintuitively) slower than splay trees.
Larger databases that require order queries also seem to be best served
by splay trees, followed by cache-oblivious \mbox{B-trees} and cache-aware
\mbox{B-trees}.

\section*{Suggestions for further research}
The dictionary problem is well-explored, including structures for efficient
in particular models of computation and practical libraries for real computers.
The space-time limitations of this thesis unfortunately did not permit us to
completely survey all of them.

We believe an benchmarking various non-traditional self-adjusting structures,
like tango trees from \cite{tango} or multi-splay trees from
\cite{multisplay-trees}, could give interesting results. Splay trees are very
performant on practical access patterns, but their worst-case $\O(N)$ behavior
limits their usefulness in time-sensitive systems. Some self-adjusting
structures provide better worst-case access times (e.g.\ $\O(\log N)$ in
multi-splay trees. On the other hand, the implementation of splay trees
are considerably simpler, so the complexity of operations might outweight
time saved by fewer cache misses in the worst case.

Several data structures based on tries have also been proposed for the
dictionary problem. One example are \emph{y-fast tries} \cite{y-fast},
which allow \textsc{Find}s in time $\O(1)$ and updates and predecessor/successor
queries in time $\O(\log\log |U|)$, where $U$ is the key universe. They require
$\O(N \log\log |U|)$ memory.
The RAM model is required by y-fast tries -- they assume that RAM operations
can be performed on keys in constant time.
The original motivation for developing y-fast tries was lowering the memory
requirements of van Emde Boas queues ($\O(|U|)$) while keeping the same time
for operations.
% TODO: Try to find some practical numbers.

\emph{Judy arrays} are highly optimized dictionaries developed at IBM
\cite{judy-shop-manual, judy-patent}.
IBM sources claim that Judy arrays are very efficient in practice.
The data structure is based on sparse 256-ary tries with implementation tricks
that compress the representation.
Unfortunately, we found the existing documentation highly complex and sometimes
incomplete. On the other hand, an open-source implementation is available
at \url{http://judy.sourceforge.net/}, so it should be possible to independently
analyze the data structure and possibly to improve upon it.


%%% Seznam použité literatury
% Seznam použité literatury je zpracován podle platných standardů. Povinnou citační
% normou pro bakalářskou práci je ISO 690. Jména časopisů lze uvádět zkráceně, ale jen
% v kodifikované podobě. Všechny použité zdroje a prameny musí být řádně citovány.

% TODO: fix citations

\def\bibname{Bibliography}
\begin{thebibliography}{99}
\addcontentsline{toc}{chapter}{\bibname}

\bibitem{ailamaki05}
	{\sc Ailamaki}, A.
	\emph{Database architectures for new hardware}.
	In \textit{ICDE '05: Proceedings of the 21st International
	Conference on Data Engineering (ICDE'05)}.
	IEEE Computer Society, Washington, DC, USA, 1148.

\bibitem{demaine00}
	% TODO: better citation
	{\sc Bender}, {\sc Demaine}, {\sc Farach-Colton}.
	\emph{Cache-Oblivious B-Trees}.
	Proceedings of the 41st Annual Symposium on
	Foundations of Computer Science, 2000.

\bibitem{brodal01}
	{\sc Brodal}, {\sc Fagenberg}, {\sc Jacob}.
	\emph{Cache Oblivious Search Trees via Binary Trees of Small Height}.
	BRICS Series RS-01-36,
	ISSN 0909-0878
	% TODO: URL

\bibitem{splay}
	{\sc Sleator}, Daniel Domonic, {\sc Tarjan}, Robert Endre.
	\emph{Self-Adjusting Binary Search Trees}.
	Journal of the ACM, Volume 32 Issue 3, July 1985, pages 652--686

\bibitem{willard92}
	{\sc Willard}, Dan E.
	\emph{A density control algorithm for doing insertions and deletions
	in a sequentially ordered file in a good worst-case time}

\bibitem{ksplay-sherk}
	{\sc Sherk}, Murray.
	\emph{Self-Adjusting $k$-ary Search Trees}

\bibitem{ksplay-nonopt}
	{\sc McClurkin}, D.J., {\sc Gaorgakopoulos}, G.F.
	\emph{Sphendamn\oe: A Proof that $k$-Splay Fails to Achieve $\log_k N$
	Behaviour}.

%\bibitem{lamport94}
%  {\sc Lamport,} Leslie.
%  \emph{\LaTeX: A Document Preparation System}.
%  2. vydání.
%  Massachusetts: Addison Wesley, 1994.
%  ISBN 0-201-52983-1.

\end{thebibliography}


%%% Tabulky v bakalářské práci, existují-li.
\chapwithtoc{List of Tables}

%%% Použité zkratky v bakalářské práci, existují-li, včetně jejich vysvětlení.
\chapwithtoc{List of Abbreviations}

%%% Přílohy k bakalářské práci, existují-li (různé dodatky jako výpisy programů,
%%% diagramy apod.). Každá příloha musí být alespoň jednou odkazována z vlastního
%%% textu práce. Přílohy se číslují.
\chapwithtoc{Attachments}

\openright
\end{document}
